% 2026-11-23 Michelman 424--427 頁（PDF 頁 28--31）與註 72--77，主對話以頁圖核對 426、427 頁之註號與引文；
% 424、425 頁依文字層。Post 432 頁與註 12--14 自 JSTOR 文字層讀過。Brennan 兩文與 Michelman〈Super Liberal〉均無本機藏本，一律轉引 Michelman 之註。
% Churchill 之語：Michelman 未註出處，正文不另寫年份。
\chapter{Brennan 之民主觀}
\label{ch:50_brennan}

\begin{qaquestion}
Michelman 第 VII 部（424--427 頁）題為「Brennan on Democracy」，Michelman\ix{persons}{michelman@Michelman, Frank I.} 在此回到 Post\ix{persons}{post@Post, Robert}，再回到 Brennan\ix{persons}{brennan@Brennan, William J., Jr.} 大法官。
\begin{enumerate}
\tightlist
\item Michelman 如何修正 Post 的直覺，使自治不以人人有權決定法律內容為前提（424--425 頁）？
\item 425--426 頁所設的兩項條件各是什麼？
\item Michelman 認為 Brennan 承認哪些非合意機關，位階如何（426 頁）？
\item Michelman 引 Brennan 的〈In Defense of Dissents〉（1985 年）與〈Reason, Passion〉（1987 年）的哪些話，用以說明什麼（426--427 頁，註 73--77）？
\item 結尾「To press your views upon ruling authorities is not yet to rule」與 Churchill 式的語句，表達 Michelman 對 Brennan 民主觀的什麼評價（427 頁）？
\end{enumerate}
\end{qaquestion}

\section{對 Post 直覺的修正}
\label{sec:50_brennan-1}

424 頁開頭，Michelman 回想前文（註 72 指示原文第 III 部 B 節）所述 Post 的直覺：若有某種實質內容是法律須含有才算正確，則就該範圍而言，「it is not rightfully up to any of us to say what the laws in substance shall be」〔法律的實質內容應為何，並不正當地由我們任何人說了算〕，因而我們不能自治，至少不能正當地自治。\footcite[424 \& n.72]{michelman1998}這項直覺在 Michelman 看來可以理解，但「may nevertheless be mistaken」〔仍可能是錯的〕：人可以自治，而不必享有、也不必自認享有，使法律成為她所決定之內容的道德權利。\footcite[424]{michelman1998}

Michelman 舉的例子是一個人自願遵守一套根本法，「out of respect for」〔出於尊重〕\ix{terms}{04@\zhlidxcat{正當性}!出於尊重}，因為她反省後認為這些法律公正、公平，或在別的意義上如其所應然。這個人可說處於「a state of self-government, in as full a sense as we can reasonably hope to find in politics」〔我們在政治中所能合理期待找到的最完整意義上的自治狀態〕。\footcite[424--425]{michelman1998}Post 在 432 頁引的就是這句話，並稱不清楚 Michelman 本人是否接受這項立場。\footcite[432 \& n.12]{post1998}

425 頁接著問，是根本法的什麼屬性使人能出於尊重而遵守。回應式民主理論（responsive-democracy theory，回應式民主\ix{terms}{01@\zhlidxcat{民主理論}!回應式民主}）可被解為：她只能尊重那些她能以公共論述的貢獻直接影響其內容的法律，縱使影響只有一小部分。但這種解釋對相信有程序獨立之正確標準的人不能成立：對這樣的人，使根本法值得尊重的屬性必須是她有理由相信那些法律正確，保有一絲控制或影響力不能充當這種屬性；Michelman 說兩者「not at all the same thing」〔完全不是一回事〕，除非她自信對道德判斷錯誤有格外的免疫。\footcite[425]{michelman1998}

\section{兩項條件}
\label{sec:50_brennan-2}

425 頁在這個基礎上說「Now complicate the situation」〔現在把情況弄複雜些〕，隨即設定一組信念，再列兩項條件。表 \ref{tab:50_brennan-conditions} 逐句錄出原文，並附 426 頁開頭的結語兩句。

\begin{longtable}{@{}L{0.07\linewidth}L{0.50\linewidth}L{0.35\linewidth}@{}}
\caption{Michelman 425--426 頁兩項條件之原文與中譯。}\label{tab:50_brennan-conditions}\\
\toprule
頁次 & 原文 & 中譯 \\
\midrule
\endfirsthead
\toprule
頁次 & 原文 & 中譯 \\
\midrule
\endhead
425 & Suppose you believe that there are procedure-independent standards of rightness for fundamental laws in this country now. & 假設你相信，這個國家現在對根本法存在著程序獨立之正確標準。 \\
425 & And suppose you think these standards involve commitments both to political self-determination through free public discourse and to respect for individual human dignity. & 又假設你認為這些標準包含兩項承諾：經由自由的公共論述達成政治自決，以及尊重個人的人性尊嚴。 \\
425 & You do not believe, however, that anyone judging in isolation---not yourself, not Justice Brennan---is a reliable authority as to exactly what these standards of rightness require of an interpreted set of fundamental laws in this country. & 然而你不相信任何孤立地判斷的人，不論是你自己或 Brennan 大法官，對於這些正確標準究竟要求這個國家一套經解釋的根本法有何內容，是可靠的權威。 \\
425 & In fact, you believe that the most reliable judge must be someone who is constantly exposed to the full blast of the sundry opinions on such questions, freely and uninhibitedly produced by assorted members of society out of their diverse experiences, positions, and reflections. & 其實你相信，最可靠的判斷者必須是不斷暴露於各種意見的全面衝擊之人，那些意見關於這類問題，由社會各成員依其不同的經驗、位置與反省，自由而無拘束地產生。 \\
425 & If you hold this set of beliefs, then the attributes in a set of fundamental laws that would command your respect would be as follows: & 若你持有這組信念，則使一套根本法得到你尊重的屬性如下： \\
425 & First, it would have to appear to you that a maximum feasible effort is constantly being made to get the fundamental laws right. This, I have argued, includes getting their major interpretations right. Quite conceivably, this maximum feasible effort could (as Dworkin contends) include provision for a tribunal whose special business and concern it is to decide the interpretations. & 第一，在你看來必須顯示，有人持續盡最大可行的努力使根本法正確。我已論證，這包括使其主要解釋正確。這項最大可行的努力很可能（如 Dworkin 所主張）包括設置一個以決定解釋為專職與專務的審判機構。 \\
425 & Second, this maximum feasible effort to get the fundamental laws right would have to include arrangements for exposing the empowered fundamental-law deciders to the full blast of sundry opinions and interest-articulations in society, including everyone's opinions and articulations of interests. & 第二，這項使根本法正確的最大可行的努力，須包括若干安排，使受授權決定根本法的人暴露於社會上各種意見與利益表達的全面衝擊，包括每個人的意見與利益表達。 \\
425 & Note that what now matters to you about having your own opinions and interest-articulations registered is not some bit of leverage you exercise over what gets decided, but rather the presumptive epistemic value of your contributions to the debate, a value which you will not self-respectingly suppose less than equal to that of others. & 請注意，你現在在意自己的意見與利益表達得以登記，所重視者為你對辯論之貢獻的推定認識價值，你自重之人不會認為這價值低於他人；你對所作決定施加的一點槓桿力，並不在其內。 \\
426 & Suppose you hold this set of beliefs, including the belief that the preceding two conditions for the production of fundamental laws (and their interpretations) are currently satisfied. You, then, can abide by the fundamental laws (including major interpretations) willingly, out of respect for them. & 假設你持有這組信念，包括相信上述產出根本法（及其解釋）的兩項條件目前已獲滿足。那麼你便能自願遵守根本法（包括主要解釋），出於對它們的尊重。 \\
426 & Are you not then freely governing yourself, at least so far as politics are concerned? & 那麼你在政治事務上，不就是自由地治理自己了嗎？ \\
\bottomrule
\end{longtable}

Post 在 432 頁以區塊引文錄入自「the attributes in a set of fundamental laws that would command your respect」起，經兩項條件，至 426 頁「Are you not then freely governing yourself」止的文字，並稱這兩項條件為「Michelman's model」〔Michelman 的模型〕。\footcite[432 \& nn.12--14]{post1998}Post 對它的批評見第~\ref{ch:60_post} 部。

第二項條件所用「full blast of sundry opinions」〔各種意見之全面衝擊〕\ix{terms}{04@\zhlidxcat{正當性}!各種意見之全面衝擊}，與 423 頁論候選屬性的用語相同。條件末句說登記自己意見的意義在於「presumptive epistemic value」〔推定的認識價值〕\ix{terms}{04@\zhlidxcat{正當性}!認識價值}，這價值是「you will not self-respectingly suppose less than equal to that of others」〔你自重之人不會認為低於他人者〕，兩個用語各對應 423 頁所說認識考量與尊嚴考量。\footcite[423, 425]{michelman1998}第一項條件提到專司解釋的審判機構，註 60 則在 419 頁聲明權利清單的設定不涉適用或執行它的機關安排。\footcite[419 n.60, 425]{michelman1998}

\section{Brennan 所列之非合意機關}
\label{sec:50_brennan-3}

Michelman 在 426 頁轉向 Brennan，寫道 Brennan 「may well have believed something along these lines」〔很可能持有與此相近的想法〕，並無疑地相信：憲政主義既預設程序獨立之正確標準\ix{terms}{04@\zhlidxcat{正當性}!程序獨立之正確標準}，又要求有某個或某些「non-consensual organ or organs」〔非合意機關〕充當社會權威，在具體議題與爭端中決定這些標準的要求。Michelman 說 Brennan 會列入這類機關者，「in due order of rank」〔依位階順序〕如表 \ref{tab:50_brennan-organs}。\footcite[426]{michelman1998}

\begin{longtable}{@{}L{0.09\linewidth}L{0.45\linewidth}L{0.38\linewidth}@{}}
\caption{Michelman 426 頁所述 Brennan 會列入之非合意機關及其位階。}\label{tab:50_brennan-organs}\\
\toprule
位階 & 原文 & 中譯 \\
\midrule
\endfirsthead
\toprule
位階 & 原文 & 中譯 \\
\midrule
\endhead
1 & Congress and state legislatures (or conventions) enacting constitutional amendments & 制定憲法修正的國會與各州議會（或制憲會議） \\
2 & conceivably an Article V constitutional convention & 設想中的依（美國聯邦憲法）第 5 條召開的制憲會議 \\
3 & constitutional courts & 憲法法院（原文未界定其範圍） \\
4 & national and state legislatures enacting ordinary legislation & 制定一般立法的國會與各州議會 \\
\bottomrule
\end{longtable}

位階一欄是 Michelman 對 Brennan 想法的描述，他沒有為這份排列引 Brennan 自己的話。依 Michelman，Brennan 另相信：自治只存在於人民出於尊重而遵守這些機關的非合意決定之時，而要使人民的尊重得到保證，除了這些機關暴露於社會上各種意見的全面衝擊，沒有任何事物足以為所需的尊重提供根據，而這種暴露須由機關自己的基本立法與解釋行為來保證。\footcite[426]{michelman1998}

Michelman 因此說，Brennan 的司法生涯反映對經由民主達成自治的承諾，他的憲法法理的每一重要面向似乎都與這組承諾與信念相合。Brennan 無疑認為他的法院被賦予權威與責任，為國家解釋一套程序獨立的正確、正義與民主標準，並認為法院獲有這項授權，僅限於法院行使權力時以保護與擴大每個人向法院及其他社會權威民主地表達意見的權利與機會為目標的期間，Michelman 說這「enormously to his credit」〔極為可貴〕。\footcite[426]{michelman1998}

註 73 稱 Brennan 法理的一項醒目特徵可說是其象徵（emblematic）：Brennan 在 NAACP v. Button\ix{cases}{02@美國!01@聯邦最高法院!1963NAACP v. Button（1963）@NAACP v. Button（1963）}, 371 U.S. 415 (1963) 為法院撰寫的意見，與在 Walker v. Birmingham\ix{cases}{02@美國!01@聯邦最高法院!1967Walker v. Birmingham（1967）@Walker v. Birmingham（1967）}, 388 U.S. 307, 346--49 (1967) 的不同意見，堅持法院本身是政治的正當場所，並且是政治鼓動（political agitation）的正當場所。註 73 另指 Brennan 一項極受矚目而高度爭議的主張：法庭，即使是上訴法庭，是法官與訴訟當事人熱烈相遇之處，即「one human being and another」〔一人與另一人〕之間的相遇，司法判斷因此是人的判斷，由法官設法同情地掌握當事人對案件的經驗而形成。註文對前者引 Michelman〈Super Liberal〉（註 22；1991 年）1268--83、1269--74 頁，對後者引 Brennan〈Reason, Passion, and 'The Progress of the Law'〉，Rec. Ass'n B. City N.Y. 948, 966 (1987)。\footcite[426 n.73]{michelman1998}

\section{Brennan 本人之文字}
\label{sec:50_brennan-4}

426 頁末到 427 頁，Michelman 以 Brennan 的〈Defense of Dissents〉及〈Reason, Passion〉為據。Michelman 正文稱前者為 1985 年，註 23 載明刊本為 37 Hastings L.J. 427 (1986)，註 74--76 的頁碼依該刊本。\footcite[426--427 \& nn.23, 74--76]{michelman1998}表 \ref{tab:50_brennan-quotes} 依 427 頁與註 74--77 列出 Michelman 所引各句。

\begin{longtable}{@{}L{0.46\linewidth}L{0.32\linewidth}L{0.14\linewidth}@{}}
\caption{Michelman 427 頁所引 Brennan 之語。}\label{tab:50_brennan-quotes}\\
\toprule
Brennan 之語 & 中譯 & 註 \\
\midrule
\endfirsthead
\toprule
Brennan 之語 & 中譯 & 註 \\
\midrule
\endhead
not only allow, [they] encourage debate & 不僅容許，也鼓勵辯論 & 74：437 頁 \\
do not shut down communication as soon as a decision is reached & 作成決定時不立即關閉溝通 & 74：437 頁 \\
our right to continue to challenge the wisdom of [any] result [is] accepted by those who disagree with us & 我們繼續質疑〔任何〕結果之智慧的權利，為與我們意見不同者所接受 & 75：同上，437 頁 \\
the ones \ldots{} that seek to sow seeds for future harvest & 那些⋯⋯尋求為將來的收成播種者（指最偉大的不同意見） & 76：430--431 頁 \\
it is only as each generation brings to bear its experience and understanding, its passion and reason & 唯有每一世代發揮其經驗與理解、其熱情與理性 & 77：962 頁 \\
hope for progress in the law & 法律進步的希望 & 77：962 頁 \\
\bottomrule
\end{longtable}

Brennan 在〈Defense of Dissents〉否認法官在高等法院中表現出的歧見會破壞公眾對司法解釋權威的接受；依 Michelman 轉述，美國人能知道自己是自由的人民，正因為且僅限於其立法者（法官包括在內）鼓勵辯論；美國人作為守法者，只在對繼續質疑結果的權利為不同意者所接受有信心時，才接受決策機關的結論。\footcite[426--427 \& nn.74--75]{michelman1998}Michelman 據此說，在 Brennan 看來最偉大的不同意見是預言性的。\footcite[427 \& n.76]{michelman1998}Michelman 又引 Brennan 在司法生涯後期所說：唯有每一世代發揮其經驗、理解、熱情與理性，法律才有進步的希望。註 77 另參見 Michelman〈Super Liberal〉（註 22）1271--72 頁。\footcite[427 \& n.77]{michelman1998}

表中前三列所述，即決策機關作成決定後不關閉溝通、容許異議者繼續質疑，與 426 頁所說「exposure, guaranteed by their own fundamental legislative and interpretative acts」〔暴露，由其自己的基本立法與解釋行為所保證〕相合。

\section{結尾}
\label{sec:50_brennan-5}

427 頁最後一段先以「If that be Brennanism, though, it still falls short of everyone's self-government」〔然而，若那就是 Brennan 主義，它仍不足以成為所有人之自治〕收束，接著寫：「To press your views upon ruling authorities is not yet to rule. To find the laws deserving of your respect is not yet to decide the laws.」〔向統治機關施加你的意見，尚不是統治。認為法律值得你尊重，尚不是決定法律。〕\footcite[427]{michelman1998}前一句的「press your views」指向決定機關表達意見，後一句的「find the laws deserving of your respect」指出於尊重而遵守；兩句所稱尚未做到的，是決定法律內容。420 頁已說明根本法的內容不能由民主程序決定，427 頁的限定與之一致。\footcite[420, 426--427]{michelman1998}

接下來是 Churchill\ix{persons}{churchill@Churchill, Winston S.} 式的語句：「Maybe, indeed, Brennanism is the worst form of democracy except all the others that have been tried from time to time.」〔或許，Brennan 主義確是最壞的民主形式，但歷來試過的其他形式除外。〕句式取自 Churchill 論民主為最壞的政府形式、但已試過的其他形式除外之說，Michelman 在文中未註明出處。他隨即追問：「And all the others that have yet to be tried?」〔至於尚未試過的其他形式呢？〕並表示不以這個結論為歡迎，寧可結束於：Brennan 給我們一項「fair test」〔公平的測試〕，提供了「the best exhibition of his version of democracy that we and our posterity are ever likely to see」〔我們與子孫可能見到的、他那種民主觀的最佳展示〕。\footcite[427]{michelman1998}

\section{課前可提的問題}
\label{sec:50_brennan-6}

\begin{enumerate}
\tightlist
\item 425 頁第二項條件以推定的認識價值界定登記意見的意義，排除對決定結果的槓桿力。公民的意見未被採納時，她是否仍可依這項條件「out of respect」遵守？
\item 424 頁「democratic pursuit of democratic justice」的首要要求（決定機關持續受公眾意見的壓力），與 425--426 頁的兩項條件，各由哪些要素構成？兩者是否重疊？
\item Post 在 432 頁以兩項條件為批評對象（第~\ref{ch:60_post} 部）。讀 Michelman 原文，426 頁以「may well have believed」把這些條件歸為 Brennan 可能持有的信念，Post 的批評針對的是 Michelman 的主張，抑或 Michelman 對 Brennan 的描述？
\end{enumerate}
